% --- INSTITUTIONAL TEMPLATE FED v7.4.9 ---
% Estándar de Reporteo APA 7 (Babel-Free)
\documentclass[12pt,spanish]{article}

\usepackage{graphicx}
\usepackage{geometry}
\usepackage{indentfirst} 
\usepackage{caption}
\usepackage{floatrow} 

% Forzar estilos de caption antes de cualquier carga
\floatsetup[figure]{capposition=top}
\floatsetup[table]{capposition=top}

\usepackage{booktabs}
\usepackage{float}
\usepackage{longtable}
\usepackage{threeparttable}
\usepackage{array}
\usepackage{calc}
\usepackage{amsmath, amssymb}
\usepackage{fancyhdr}
\usepackage{mathptmx} 
\usepackage[T1]{fontenc}
\usepackage{setspace} 
\usepackage{ragged2e} 
\usepackage{titlesec}
\usepackage{url}
\def\UrlBreaks{\do\/\do-\do\.\do\_\do\?\do\=\do\&\do\%}
\usepackage[hidelinks,breaklinks]{hyperref}

% --- Pandoc Compatibility Block (Iron Clad) ---
\providecommand{\tightlist}{%
  \setlength{\itemsep}{0pt}\setlength{\parskip}{0pt}}

% Compatibilidad Pandoc 3.x
\ifdefined\pandocbounded\else
  \newcommand{\pandocbounded}[1]{#1}
\fi

\makeatletter
\@ifundefined{paragraph}{\newcounter{paragraph}}{}
\@ifundefined{subparagraph}{\newcounter{subparagraph}}{}
\makeatother

% Solución al error 'No counter none defined' en longtable
\makeatletter
\newcounter{none}
\@ifpackageloaded{caption}{
  \DeclareCaptionType{none}
  \captionsetup[none]{name=Tabla}
}{}
\makeatother

\hypersetup{
  colorlinks=true,
  linkcolor=blue,
  filecolor=blue,
  citecolor=blue,
  urlcolor=blue
}

% --- Macros Institucionales ---
\newcommand{\fuente}[1]{
  \vspace{4pt}
  \begin{flushleft}
    \small
    \textit{Nota.} #1
  \end{flushleft}
}

% --- Localización Manual Determinista (Sin Babel) ---
\AtBeginDocument{
  \renewcommand{\tablename}{Tabla}
  \renewcommand{\figurename}{Figura}
  \renewcommand{\contentsname}{Índice}
  \renewcommand{\listtablename}{Índice de Tablas}
  \renewcommand{\listfigurename}{Índice de Figuras}
  \renewcommand{\abstractname}{Resumen}
  \renewcommand{\refname}{Referencias}
  
  \RaggedRight
  \setlength{\parindent}{1.27cm}
  \setlength{\parskip}{0pt}
}

\DeclareCaptionLabelSeparator{newline}{\\ \vspace{2pt}}
\captionsetup{
  position=top,
  justification=RaggedRight,
  singlelinecheck=false,
  labelfont=bf,
  textfont=it,
  labelsep=newline,
  font=small,
  skip=10pt
}
\captionsetup[figure]{name=Figura}
\captionsetup[table]{name=Tabla}

% Márgenes APA 7
\geometry{
  a4paper,
  margin=25.4mm,
  headheight=15mm
}

% --- Encabezados ---
\pagestyle{fancy}
\fancyhf{} 
\fancyhead[L]{\includegraphics[height=1.2cm]{/home/erick-fcs/Capital_Workstation/capital-workstation-libs/src/ecs_quantitative/reporting/assets/logos/logo_unl.png}}
\fancyhead[R]{\includegraphics[height=1.2cm]{/home/erick-fcs/Capital_Workstation/capital-workstation-libs/src/ecs_quantitative/reporting/assets/logos/logo_economia.jpg}}
\renewcommand{\headrulewidth}{0pt}
\fancyfoot[R]{\thepage}

\fancypagestyle{plain}{
  \fancyhf{}
  \fancyhead[L]{\includegraphics[height=1.2cm]{/home/erick-fcs/Capital_Workstation/capital-workstation-libs/src/ecs_quantitative/reporting/assets/logos/logo_unl.png}}
  \fancyhead[R]{\includegraphics[height=1.2cm]{/home/erick-fcs/Capital_Workstation/capital-workstation-libs/src/ecs_quantitative/reporting/assets/logos/logo_economia.jpg}}
  \fancyfoot[R]{\thepage}
  \renewcommand{\headrulewidth}{0pt}
}

% --- Títulos ---
\titleformat{\section}{\normalfont\normalsize\bfseries\centering}{\thesection}{1em}{}
\titleformat{\subsection}{\normalfont\normalsize\bfseries}{\thesubsection}{1em}{}
\titleformat{\subsubsection}{\normalfont\normalsize\bfseries\itshape}{\thesubsubsection}{1em}{}
\titlespacing*{\section}{0pt}{12pt}{6pt}
\titlespacing*{\subsection}{0pt}{12pt}{6pt}
\titlespacing*{\subsubsection}{0pt}{12pt}{6pt}

\newenvironment{referenciasAPA}{
  \setlength{\parindent}{-1.27cm}
  \setlength{\leftskip}{1.27cm}
  \setlength{\parskip}{6pt}
}{\par}

\begin{document}

\begin{center}
    {\bfseries \Large UNIVERSIDAD NACIONAL DE LOJA}\\[0.3cm]
    {\bfseries \large FACULTAD JURÍDICA, SOCIAL Y
ADMINISTRATIVA}\\[0.2cm]
    {\bfseries CARRERA DE ECONOMÍA}\\[1.5cm]
    
    {\bfseries \large POLÍTICA ECONÓMICA}\\[1cm]
    
    {\bfseries \Large Actividad AA: Aplicación del Marco
Conceptual}\\[1.5cm]
\end{center}

\noindent {\bfseries Estudiante:} Erick Fabricio Condoy Seraquive \\
\noindent {\bfseries Docente:} Econ. Maria Gabriela Moreno Hurtado \\
\noindent {\bfseries Fecha:} 19 de abril de 2026 \\

\vspace{1cm}

\doublespacing
\setlength{\parindent}{1.27cm}

\section{Identificación y
Clasificación}\label{identificaciuxf3n-y-clasificaciuxf3n}

Para el presente análisis, se ha seleccionado el Servicio Ampliado del
Fondo (SAF/EFF) alcanzado entre el Gobierno de Ecuador y el Fondo
Monetario Internacional (FMI) en mayo de 2024, por un monto de US\$
4.000 millones y una duración de 48 meses según la publicación oficial
del FMI (2024).

\subsection{Taxonomía de la Política bajo la visión de Cuadrado
Roura}\label{taxonomuxeda-de-la-poluxedtica-bajo-la-visiuxf3n-de-cuadrado-roura}

La clasificación tipológica de esta política, fundamentada en los
criterios taxonómicos clásicos, se define de la siguiente manera:

\begin{itemize}
\tightlist
\item
  Orientación Básica: Se clasifica como una política de proceso y de
  ordenación. Según Cuadrado Roura (2014) en su obra sobre Política
  Económica, las políticas de proceso actúan sobre los desequilibrios
  coyunturales, mientras que las de ordenación buscan establecer o
  modificar el marco de actuación de la economía.
\item
  Carácter de los Instrumentos: Es predominantemente una política
  cuantitativa combinada con políticas cualitativas que conllevan
  reformas. Las variaciones en el IVA fijado en el 15\% son ajustes
  cuantitativos, mientras que la focalización de subsidios representa un
  cambio cualitativo estructural.
\item
  Nivel de Actuación: Taxonómicamente es una política macroeconómica. Su
  núcleo de atención son los agregados macroeconómicos para lograr la
  estabilidad bi-unívoca entre sostenibilidad fiscal y balanza de pagos.
\item
  Dimensión Temporal: Se define como una política de mediano plazo, con
  un horizonte de ejecución de 48 meses.
\end{itemize}

\section{Fines, Objetivos e
Instrumentos}\label{fines-objetivos-e-instrumentos}

Siguiendo la lógica de la ``jerarquía de la política económica'',
distinguimos los siguientes niveles en el Acuerdo SAF 2024:

\subsection{Fines Generales definidos como Ejes
Estratégicos}\label{fines-generales-definidos-como-ejes-estratuxe9gicos}

Los fines que motivan esta política son la estabilidad institucional y
la sostenibilidad del sistema monetario con el fin de blindar la
dolarización. Estos representan propósitos de carácter genérico que solo
tienen sentido en la medida en que se traducen en objetivos concretos.

\subsection{Objetivos Económicos y
Sociales}\label{objetivos-econuxf3micos-y-sociales}

\begin{enumerate}
\def\labelenumi{\arabic{enumi}.}
\tightlist
\item
  Consolidación Fiscal: Reducción del déficit primario.
\item
  Protección Social: Establecimiento de pisos de gasto social.
\item
  Estabilidad Externa: Acumulación de reservas internacionales.
\end{enumerate}

\subsection{Instrumentos Específicos}\label{instrumentos-especuxedficos}

\begin{itemize}
\tightlist
\item
  Tributarios: Incremento del Impuesto al Valor Agregado (IVA) al 15\%
  (Primicias, 2024).
\item
  Gasto Público: Focalización de los subsidios a los combustibles y
  optimización del gasto.
\end{itemize}

\section{Relación
Instrumentos-Objetivos}\label{relaciuxf3n-instrumentos-objetivos}

La siguiente matriz presenta el mapeo de sinergias y conflictos entre
los principales instrumentos del Acuerdo SAF 2024 y los objetivos de la
política económica analizada.

{

\begin{longtable}[]{@{}
  >{\raggedright\arraybackslash}p{(\linewidth - 10\tabcolsep) * \real{0.3500}}
  >{\centering\arraybackslash}p{(\linewidth - 10\tabcolsep) * \real{0.1300}}
  >{\centering\arraybackslash}p{(\linewidth - 10\tabcolsep) * \real{0.1300}}
  >{\centering\arraybackslash}p{(\linewidth - 10\tabcolsep) * \real{0.1300}}
  >{\centering\arraybackslash}p{(\linewidth - 10\tabcolsep) * \real{0.1300}}
  >{\centering\arraybackslash}p{(\linewidth - 10\tabcolsep) * \real{0.1300}}@{}}

\caption{\label{tbl-matriz}Matriz de Relación entre Instrumentos y
Objetivos de Política Económica (2024)}

\tabularnewline

\toprule\noalign{}
\begin{minipage}[b]{\linewidth}\raggedright
Instrumento
\end{minipage} & \begin{minipage}[b]{\linewidth}\centering
Crecimiento
\end{minipage} & \begin{minipage}[b]{\linewidth}\centering
Empleo
\end{minipage} & \begin{minipage}[b]{\linewidth}\centering
Fiscal
\end{minipage} & \begin{minipage}[b]{\linewidth}\centering
Equidad
\end{minipage} & \begin{minipage}[b]{\linewidth}\centering
Exterior
\end{minipage} \\
\midrule\noalign{}
\endhead
\bottomrule\noalign{}
\endlastfoot
Alza de IVA al 15\% & - & & + & & + \\
Focalización Subsidios & & & + & + & \\
Pisos de Gasto Social & & & & + & \\

\end{longtable}

}

\fuente{Elaboración propia con base en el marco teórico de Cuadrado Roura (2014), el Acuerdo SAF 2024 (IMF, 2024) y datos de seguimiento de Primicias (2024).}

\subsection{Análisis de Coherencia sobre el Objetivo de
Empleo}\label{anuxe1lisis-de-coherencia-sobre-el-objetivo-de-empleo}

Es pertinente observar que los instrumentos analizados no muestran un
efecto positivo directo sobre el indicador de empleo. Esta neutralidad
es económicamente coherente con la naturaleza del Acuerdo SAF 2024, el
cual es un programa de estabilización y consolidación fiscal, no de
estímulo directo a la demanda agregada. La teoría económica advierte que
los ajustes impositivos (IVA) y la reducción de subsidios tienen efectos
contractivos en el corto plazo, priorizando la solvencia del Estado
sobre la expansión inmediata del mercado laboral.

\subsection{Conclusión del
Análisis}\label{conclusiuxf3n-del-anuxe1lisis}

La política económica implementada muestra una fuerte prioridad hacia la
sostenibilidad fiscal. El cumplimiento de las metas del FMI se presenta
como la vía para garantizar la solvencia del Estado, asumiendo costos
directos sobre el consumo pero buscando compensar a los sectores
vulnerables mediante transferencias monetarias condicionadas, tal como
se detalla en el comunicado oficial 24/197 del FMI.

\newpage

\section{Referencias}\label{referencias}

\begin{referenciasAPA}
Cuadrado Roura, J. R. (2014). *Política Económica: Elaboración, objetivos e instrumentos*. Quinta edición. Madrid: McGraw-Hill.

International Monetary Fund. (2024, 31 de mayo). *IMF Executive Board Approves 48-Month US\$4 Billion EFF Arrangement for Ecuador*. Press Release No. 24/197. \url{https://www.imf.org/en/news/articles/2024/05/31/pr-24197-ecuador-imf-board-approves-48-month-us-4-billion-eff-arrangement}

Primicias. (2024). *FMI y Ecuador llegan a un acuerdo técnico para nuevo crédito por USD 4.000 millones*. \url{https://www.primicias.ec/noticias/economia/fmi-acuerdo-tecnico-daniel-noboa-ecuador/}
\end{referenciasAPA}

\end{document}
